% bd-screen.tex

% driver file bd-screen.tex to produce the Boxes and Diamonds textbook on
% with same type block as in printed version, but but with on-screen
% features (color, links, etc).

% We use the memoir class for maximal flexibility of layout, but any
% class will do

\documentclass[11pt,openany]{memoir}

\RequirePackage{amsthm}
\RequirePackage{xcolor}
\RequirePackage{mdframed}
\RequirePackage[full]{textcomp}
\RequirePackage{gitinfo2}

\RequirePackage[pdftex,breaklinks,bookmarks,bookmarksopen,bookmarksopenlevel=1,colorlinks,urlcolor=dark-gray,linkcolor=reflex-blue]{hyperref}

%\usepackage[nonumberlist,toc,style=index]{glossaries}
%\makeglossaries


% let's set the whole thing in Baskervald X, with Universalis ADF
% Standard for sans-serif, and spread the lines a bit to make the text
% more readable

\usepackage[sfdefault]{universalis}
\usepackage[osf]{Baskervaldx} % oldstyle figures
\usepackage[bigdelims,baskervaldx]{newtxmath} 
\usepackage[cal=boondoxo]{mathalfa}

% Make sure we have a copyright symbol

\def\copyright{\textcircled{C}}

\def\oljobname{bd}

% next line needed for Adobe Acrobat to open the PDF.  may have
% somethign to do with transparency in the PNG graphics files?

\pdfminorversion=4

% load textpos for the cover

\usepackage[absolute,overlay]{textpos}

% set stock & paper size with narrow margins

\setstocksize{22cm}{17cm}

\settrimmedsize{\stockheight}{\stockwidth}{*}
\settrims{0pt}{0pt}

% let's calculate the line length for 65 characters in \normalfont

\setlxvchars

% set the size of the type block to golden ratio calculated width

\settypeblocksize{*}{1.05\lxvchars}{1.62}

% set spine and and edge margin

\setlrmargins{*}{*}{1}
\setulmargins{60pt}{*}{*}
\setheaderspaces{*}{*}{1}

\checkandfixthelayout

% Chapter style

\makeatletter

\definecolor{dkleadbeater}{RGB}{173, 94, 153}
\definecolor{leadbeater}{RGB}{124,67,110}
\definecolor{ltleadbeater}{RGB}{239,189,227}
\newlength{\barlength}
\makechapterstyle{leadbeater}{%
  \setlength{\afterchapskip}{40pt}
  \setlength{\beforechapskip}{50pt}
    \setlength{\midchapskip}{10pt}
  \renewcommand*{\afterchapternum}{\par\nobreak\vskip 0pt}
  \renewcommand*{\chapnamefont}{\fontsize{14pt}{0pt}\selectfont\sffamily\bfseries}
  \let\chapnumfont\chapnamefont
  \renewcommand*{\chaptitlefont}{\normalfont\fontsize{48pt}{48pt}\selectfont\bfseries\itshape\color{leadbeater}}
  \renewcommand*{\printchaptername}{%
    \chapnamefont\MakeUppercase{\@chapapp}}
  \renewcommand*{\printchaptertitle}[1]{%
    \chaptitlefont ##1\\[-\baselineskip]%
    \hspace*{-20pt}%
    \smash{\color{leadbeater}\rule{7pt}{300pt}}}
}

\renewcommand*{\partnamefont}{\fontsize{24pt}{0pt}\selectfont\bfseries\sffamily}
\renewcommand*{\partnumfont}{\fontsize{24pt}{0pt}\selectfont\bfseries\sffamily}
\renewcommand*{\parttitlefont}{\normalfont\fontsize{54pt}{54pt}\selectfont\bfseries\itshape\color{leadbeater}}
\renewcommand*{\printpartname}{%
  \partnamefont PART}

\makeatother

\chapterstyle{leadbeater}

\copypagestyle{leadbeater}{headings}

\makeoddhead{leadbeater}{\small\sffamily\color{leadbeater}\rightmark}{}
            {\color{leadbeater}\sffamily\bfseries\thepage}

\makeevenhead{leadbeater}{\small\sffamily\color{leadbeater}\leftmark}{}
            {\color{leadbeater}\sffamily\bfseries\thepage}

\usepackage[font={small,it}]{caption}


% \olpath has to point to the location of the OLP main
% directory/folder.  We're compiling from subdirectory
% courses/phil379/, so the main directory is two
% levels up.

\newcommand{\olpath}{../../}

% load all the Open Logic definitions. This will also load the
% local definitions in open-logic-sample-config.sty

\input{\olpath/sty/open-logic.sty}

% we want all the problems deferred to the very end

\input{\olpath/sty/open-logic-defer.sty}

% colors for links

\hypersetup{
colorlinks,
  allcolors=dkleadbeater,
  pdftitle={Boxes and Diamonds},
  pdfauthor = {Open Logic Project}
}

% load glossary entries

%\loadglsentries{include/glossary}

% end preamble

\input{includeonly}

\begin{document}
\raggedbottom

\thispagestyle{empty}

\pagecolor{ltleadbeater}
  \begin{textblock*}{4.5in}(.25in,.75in)%
    \begin{raggedright}
      \fontsize{26pt}{30pt}\selectfont\bfseries\sffamily%
      Boxes and Diamonds\\[5pt]
      \normalfont\fontsize{18pt}{0pt}\selectfont\bfseries\itshape%
      \rule{4.5in}{5pt}\\[5pt]
      An Open Introduction to Modal Logic

\vskip1.5cm

\hfill\includegraphics[width=.29\textwidth]{\olpath/assets/portraits/barcan-circle.pdf}
\includegraphics[width=.29\textwidth]{\olpath/assets/portraits/carnap-circle.pdf}\hfill{}

\includegraphics[width=.29\textwidth]{\olpath/assets/portraits/antonelli-circle.pdf}\hfill
\includegraphics[width=.29\textwidth]{\olpath/assets/portraits/lewis-circle.pdf}\hfill
\includegraphics[width=.29\textwidth]{\olpath/assets/portraits/prior-circle.pdf}

\hfill\includegraphics[width=.29\textwidth]{\olpath/assets/portraits/heyting-circle.pdf}
\includegraphics[width=.29\textwidth]{\olpath/assets/portraits/kripke-circle.pdf}\hfill{}

\vskip1.5cm

\includegraphics[width=1cm]{assets/cc.pdf}
\includegraphics[width=1cm]{assets/by.pdf}
\includegraphics[width=1cm]{assets/remix.pdf}
\normalfont\fontsize{16pt}{0pt}\selectfont\bfseries\sffamily%
\hfill Winter 2018 \emph{bis}
    \end{raggedright}   
\end{textblock*}
\

\clearpage
\setcounter{page}{1}
\nopagecolor




% Now load the actual text

% bd.tex
%
% driver file bd.tex to produce text

\preto\OLEndChapterHook{\IfFileExists{include/summary-\thechapter}
        {\section*{Summary}\addcontentsline{toc}{section}{Summary}
        \let\emph\textbf\input{include/summary-\thechapter}\let\emph\textit}{}}

\problemsperchapter

\allowdisplaybreaks

\frontmatter


% bastard title

\pagestyle{empty}

\vspace*{100pt}

\begin{raggedleft}

{\fontsize{24pt}{24pt}\selectfont\bfseries\sffamily%
Boxes and Diamonds}

\end{raggedleft}


\newpage

% editorial board

\vspace*{100pt}

{\bfseries\itshape The Open Logic Project}

\bigskip

\textbf{\color{leadbeater}Instigator}

\medskip

Richard Zach, \emph{University of Calgary}

\bigskip

\textbf{\color{leadbeater}Editorial Board}

\medskip

Aldo Antonelli,$^\dagger$ \emph{University of California, Davis}

Andrew Arana, \emph{Universit\'e Paris I Panth\'enon--Sorbonne}

Jeremy Avigad, \emph{Carnegie Mellon University}

Walter Dean, \emph{University of Warwick}

Gillian Russell, \emph{University of North Carolina}

Nicole Wyatt, \emph{University of Calgary}

Audrey Yap, \emph{University of Victoria}

\bigskip

\textbf{\color{leadbeater}Contributors}

\medskip

Samara Burns, \emph{University of Calgary}

Dana H\"agg, \emph{University of Calgary}

Zesen Qian, \emph{Carnegie Mellon University}

\newpage

% title

\vspace*{100pt}

\begin{raggedleft}

{\fontsize{24pt}{24pt}\selectfont\bfseries\sffamily%
Boxes and Diamonds}

\bigskip

{\fontsize{18pt}{18pt}\selectfont\bfseries\itshape An Open Introduction to Modal Logic}

\vspace{100pt}

\fontsize{14pt}{14pt}\selectfont Remixed by Richard Zach

\vfill

\textsc{Winter 2018}

\end{raggedleft}


\newpage

% credits

\noindent
The Open Logic Project would like to acknowledge the generous support
of the \href{http://arts.ucalgary.ca/}{Faculty of Arts} and
the \href{http://www.ucalgary.ca/taylorinstitute/}{Taylor Institute of
Teaching and Learning} of the University of Calgary,
and the \href{http://albertaoer.com}{Alberta Open Educational Resources
(ABOER) Initiative}, which is made possible through an investment from
the Alberta government.

\bigskip

\noindent\includegraphics[height=1cm]{assets/ucarts-color}

\medskip

\noindent\includegraphics[height=1cm]{assets/ti-color}

\medskip

\noindent\includegraphics[width=4cm]{assets/aboer-color}

\vfill

\noindent Cover illustrations by \href{http://mattleadbeater.com}{Matthew
  Leadbeater}, used under a
\href{http://creativecommons.org/licenses/by-nc/4.0/}{Creative Commons
  Attribution-NonCommercial 4.0 International License}.

\vfill

\noindent Typeset in Baskervald X and Universalis ADF Standard by
\LaTeX.

\vfill

\noindent This version of \texttt{boxes-and-diamonds} is
revision \gitAbbrevHash{} (\gitAuthorDate), with content generated
from \texttt{OpenLogicProject} revision
4dadf83
(2018-01-10).


\vfill

% oluselicense generates a license mark that a) licenses the result
% under a CC-BY licence and b) acknowledges the original source (the
% OLP).  Acknowledgment of the source is a requirement under the
% conditions of the CC-BY license used by the OLP, but you are not
% required to license the product itself under CC-BY.

\renewcommand{\ollicensefont}{\fontsize{8.5pt}{11pt}\selectfont}

\noindent
\oluselicense
% Title of this version of the OLT with link to source
{\href{https://github.com/rzach/boxes-and-diamonds}{\textit{Boxes and Diamonds}}}
% Author of this version
{\href{http://richardzach.org/}{Richard Zach}}

\newpage
\pagestyle{leadbeater}
\currentpdfbookmark{Table of Contents}{name}
\tableofcontents*



\mainmatter

\pagestyle{leadbeater}

\olimport*[normal-modal-logic/syntax-and-semantics]{syntax-and-semantics}


\olimport*[normal-modal-logic/axioms-systems]{axioms-systems}


\olimport*[normal-modal-logic/completeness]{completeness}


\olimport*[normal-modal-logic/filtrations]{filtrations}


\appendix

\olimport*[sets-functions-relations/sets]{sets}
\olimport*[sets-functions-relations/relations]{relations}
\olimport*[methods/induction]{induction}
\olimport*[methods/proofs]{proofs}


\stopproblems
\def\ifproblems#1{}

\def\captionnamefont{\sffamily\color{leadbeater}}
\def\captiontitlefont{\sffamily\color{leadbeater}}

\def\figurename{Fig.}
\setlength{\olphotowidth}{.45\textwidth}

\backmatter

\clearpage
\photocredits

\bibliographystyle{\olpath/bib/natbib-oup}

\bibliography{\olpath/bib/open-logic.bib}

\olimport*{\olpath/content/open-logic-about}




\end{document}

